\documentclass[conference]{IEEEtran}
\IEEEoverridecommandlockouts
\usepackage{cite}
\usepackage{amsmath, amssymb, amsthm}
\def\BibTeX{{\rm B\kern-.05em{\sc i\kern-.025em b}\kern-.08em
    T\kern-.1667em\lower.7ex\hbox{E}\kern-.125emX}}

\begin{document}

\section*{Mathematical Formulation of the Baseline-Reverting Spiking Diffusion Synthesizer}

To develop a generative digital twin of the cortex that accounts for stationary subject-specific biometrics, we model the degradation of transient emotional brainwaves as a continuous-time Stochastic Differential Equation (SDE). Unlike standard image diffusion which drifts toward isotropic Gaussian noise $\mathcal{N}(0, I)$, our forward process is governed by a modified Ornstein-Uhlenbeck (OU) process that drifts the active state toward $C$, the subject's stationary Individual Alpha Frequency (IAF) covariance baseline.

\subsection*{1. The Forward and Reverse-Time SDE}
Let $X_t$ be the EEG state at time $t$, $C$ be the stationary subject baseline, and $\beta(t)$ be the continuous noise schedule. The forward SDE is defined as:
\begin{equation}
    dX_t = \frac{1}{2}\beta(t)(C - X_t) dt + \sqrt{\beta(t)} dW_t
\end{equation}
Where $W_t$ is the standard Wiener process. To generate novel emotional samples, we must reverse this process. By applying Anderson's Theorem \cite{Anderson1982}, the corresponding reverse-time SDE is given by:
\begin{equation}
    d\tilde{X}_t = \left[ \frac{1}{2}\beta(t)(C - \tilde{X}_t) - \beta(t) \nabla_{x} \log p_t(\tilde{X}_t) \right] dt + \sqrt{\beta(t)} d\tilde{W}_t
\end{equation}
where $dt$ represents an infinitesimal negative time step, and $d\tilde{W}_t$ is the reverse-time Wiener process. The term $\nabla_{x} \log p_t(X_t)$ represents the marginal score function, which guides the trajectory away from the baseline $C$ and injects emotional semantics.

\subsection*{2. The Transition Kernel}
To train a network to estimate this score, we must first determine the transition probability $p_{0t}(X_t \mid X_0)$. By multiplying the forward SDE by an integrating factor, we solve the linear SDE to find the exact Gaussian transition kernel:
\begin{equation}
    p_{0t}(X_t \mid X_0) = \mathcal{N}\Big(X_t ;\; X_0 e^{-A(t)} + C(1 - e^{-A(t)}), \; (1 - e^{-2A(t)})I \Big)
\end{equation}
where $A(t) = \int_0^t \frac{1}{2}\beta(s) ds$. 

This closed-form solution elegantly captures the physical decay of the emotion. At $t=0$ (where $A(0)=0$), the mean is exactly $X_0$ (the pure emotional climax). As $t \to \infty$, $e^{-A(t)} \to 0$, and the mean converges exactly to $C$, proving that the emotion decays purely into the subject's stationary IAF biometric fingerprint.

\subsection*{3. Denoising Score Matching via Reparameterization}
Training a neural network to directly predict the marginal score $\nabla_{x} \log p_t(X_t)$ is computationally intractable because the marginal distribution requires integrating over all possible initial states $X_0$. However, the transition score $\nabla_{X_t} \log p_{0t}(X_t \mid X_0)$ is tractable. We apply Denoising Score Matching (DSM, \cite{Song2019}) to bypass this intractability. 

Taking the gradient of the log-likelihood of our Gaussian transition kernel, we obtain:
\begin{equation}
    \nabla_{X_t} \log p_{0t}(X_t \mid X_0) = -\frac{X_t - \mu_t}{\Sigma_t}
\end{equation}
We define $\sigma_t = \sqrt{1 - e^{-2A(t)}}$. By applying the reparameterization trick, $X_t = \mu_t + \sigma_t \epsilon$, where $\epsilon \sim \mathcal{N}(0, I)$, the transition score simplifies to:
\begin{equation}
    \nabla_{X_t} \log p_{0t}(X_t \mid X_0) = -\frac{\epsilon}{\sigma_t}
\end{equation}

Instead of predicting the gradient directly, we parameterize our Spiking Neural Network to predict the injected noise $\epsilon_\theta(X_t, t)$. By weighting the loss objective by $\lambda(t) = \sigma_t^2$, the variance denominator is perfectly canceled out. This yields a highly stable, baseline-agnostic Mean Squared Error (MSE) objective:
\begin{equation}
    \mathcal{L}_{DSM}(\theta) = \mathbb{E}_{t, X_0, \epsilon} \left[ \| \epsilon_\theta(X_t, t) - \epsilon \|_2^2 \right]
\end{equation}
Crucially, the subject baseline constant $C$ completely drops out of the network's loss function, allowing the SNN to optimize purely on the transient emotional delta.

\subsection*{4. Fast Sampling via Probability Flow ODE}
While the reverse SDE provides exact stochastic sampling, simulating biological action potentials requires high temporal resolution which suffers under Euler-Maruyama discretization. Alternatively, the process can be converted into a deterministic Probability Flow ODE. To mitigate numerical integration/truncation errors during fast inference, we propose utilizing second-order deterministic solvers, specifically the Heun-based preconditioned solver \cite{Karras2022}, reducing the required network evaluations for real-time BCI synthesis.

\bibliographystyle{IEEEtran}
\bibliography{refs}
\end{document}

